\paragraph{\textbf{5.7.}}
\SGARefTarget{sga3:xx:5.7}
Let \((G,T,\alpha)\) now be an arbitrary elementary \(S\)-system.
Consider the commutative diagram
\begin{center}
\begin{tikzpicture}[>=stealth,
  every node/.style={inner sep=2pt},
  every path/.style={line width=.45pt}]
  \node (top) at (0,2.9) {$\mathbb G_{m,S}$};
  \node (left) at (-3.7,0) {$\mathbb G_{m,S}$};
  \node (mid) at (0,0) {$T$};
  \node (right) at (3.7,0) {$\mathbb G_{m,S}$};
  \node (bottom) at (0,-2.9) {$\mathbb G_{m,S}$};

  \draw[->] (top) -- node[pos=.48,above left] {$\mathrm{id}$} (left);
  \draw[->] (top) -- node[left] {$\alpha^\ast$} (mid);
  \draw[->] (top) -- node[pos=.48,above right] {$2$} (right);
  \draw[->] (left) -- node[above] {$\alpha^\ast$} (mid);
  \draw[->] (mid) -- node[above] {$\alpha$} (right);
  \draw[->] (left) -- node[pos=.48,below left] {$2$} (bottom);
  \draw[->] (mid) -- node[left] {$\alpha$} (bottom);
  \draw[->] (right) -- node[pos=.48,below right] {$\mathrm{id}$} (bottom);
\end{tikzpicture}
\end{center}
Combining this diagram with that of \SGARefLink{sga3:xx:5.6}{5.6} gives a commutative diagram
\begin{center}
\begin{tikzpicture}[>=stealth,
  every node/.style={inner sep=2pt},
  every path/.style={line width=.45pt}]
  \node (top) at (0,2.9) {$\mathbb G_{m,S}$};
  \node (left) at (-4.8,0) {$t_S(\mathbb G_{m,S})$};
  \node (mid) at (0,0) {$T$};
  \node (right) at (4.8,0) {$t_P(\mathbb G_{m,S})$};
  \node (bottom) at (0,-2.9) {$\mathbb G_{m,S}$};

  \draw[->] (top) -- node[pos=.5,above left] {$\alpha_S^\ast$} (left);
  \draw[->] (top) -- node[left] {$\alpha^\ast$} (mid);
  \draw[->] (top) -- node[pos=.5,above right] {$\alpha_P^\ast$} (right);
  \draw[->] (left) -- node[above] {$\alpha^\ast\circ t_S^{-1}$} (mid);
  \draw[->] (mid) -- node[above] {$t_P\circ\alpha$} (right);
  \draw[->] (left) -- node[pos=.5,below left] {$\alpha_S$} (bottom);
  \draw[->] (mid) -- node[left] {$\alpha$} (bottom);
  \draw[->] (right) -- node[pos=.5,below right] {$\alpha_P$} (bottom);
\end{tikzpicture}
\end{center}
Using \SGARefLink{sga3:xx:4.1}{4.1}, one obtains:

\paragraph{\textbf{Proposition 5.8.}}
\SGARefTarget{sga3:xx:5.8}
\emph{Let \(S\) be a scheme and let \((G,T,\alpha)\) be an
elementary \(S\)-system.  Put \(L=\mathfrak g^\alpha\), so that
\(L^{-1}=\mathfrak g^{-\alpha}\).}

\emph{\textup{(i)} There exists a unique group morphism
\(f\colon\mathrm{SL}_L\to G\) satisfying the following equivalent
conditions:}
\[
\begin{array}{ll}
\textup{(a)}
&
f\begin{pmatrix}z&0\\0&z^{-1}\end{pmatrix}=\alpha^\ast(z),
\qquad
f\begin{pmatrix}1&X\\0&1\end{pmatrix}=\exp(X);
\\[1.1em]
\textup{(b)}
&
f\begin{pmatrix}1&X\\0&1\end{pmatrix}=\exp(X),
\qquad
f\begin{pmatrix}1&0\\Y&1\end{pmatrix}=\exp(Y);
\\[1.1em]
\textup{(c)}
&
f\begin{pmatrix}1&X\\0&1\end{pmatrix}=\exp(X),
\qquad
f\begin{pmatrix}0&X\\-X^{-1}&0\end{pmatrix}=w_\alpha(X).
\end{array}
\]

\emph{\textup{(ii)} There exists a unique group morphism
\(g\colon G\to\mathrm P_L\) satisfying}
\[
 g(t)=
 \begin{pmatrix}\alpha(t)&0\\0&1\end{pmatrix},
 \qquad
 g(\exp(X))
 =p\begin{pmatrix}1&X\\0&1\end{pmatrix}.
\]
\emph{Moreover,}
\[
 g(\exp(Y))
 =p\begin{pmatrix}1&0\\Y&1\end{pmatrix},
 \qquad
 g(w_\alpha(X))
 =p\begin{pmatrix}0&X\\-X^{-1}&0\end{pmatrix}.
\]
\emph{The morphism \(g\) is faithfully flat and quasi-compact, with}
\[
 \Ker(g)=\Ker(\alpha)=\operatorname{Centr}(G),
\]
\emph{and \(g\circ f\) is the canonical morphism
\(\mathrm{SL}_L\to\mathrm P_L\).}

Notice that condition \textup{(b)} in \textup{(i)} gives an explicit
description of the duality between
\(\mathfrak g^\alpha\) and \(\mathfrak g^{-\alpha}\).

\paragraph{\textbf{Corollary 5.9.}}
\SGARefTarget{sga3:xx:5.9}
\emph{Let \((G,T,\alpha)\) be an elementary \(S\)-system.  The
subgroups}
\[
 T\cdot U_\alpha,\qquad T\cdot U_{-\alpha},
 \qquad U_\alpha,\qquad U_{-\alpha}
\]
\emph{are closed.}

Since \(U_\alpha\) is a closed subgroup scheme of
\(T\cdot U_\alpha\), it suffices to prove the assertion for the
latter.  By Noether's theorem \textup{(\SGARefLink{sga3:iv:proposition:5:3:1}{Exposé IV, 5.3.1} and \SGARefLink{sga3:iv:lemma:6:4:1}{6.4.1})},
it is enough to show that
\[
 (T\cdot U_\alpha)/\Ker(\alpha)
\]
is a closed subgroup of \(G/\Ker(\alpha)\).  By \SGARefLink{sga3:xx:5.8}{5.8}, this reduces to
the evident fact that the subgroup of \(\mathrm P_L\) defined by
\(c=0\) is closed; one may equivalently work in \(\mathrm{GL}_L\),
by another application of Noether's theorem.  Consequently, the
exponential morphisms of \SGARefLink{sga3:xx:1.5}{Theorem 1.5}\textup{(i)} are closed
immersions.

\smallskip
\noindent\textbf{N.B.}\ The corollary also follows from the fact that
\(T\cdot U_\alpha\) and \(T\cdot U_{-\alpha}\) are ``Borel
subgroups'' of \(G\) \textup{(cf.\ \SGARefLink{sga3:xii:7.10}{Exposé XII, 7.10})}.

\paragraph{\textbf{5.10.}}
\SGARefTarget{sga3:xx:5.10}
Let \(L\) be an invertible \(\mathcal O_S\)-module, and let
\[
 \mathbb G_{m,S}\xrightarrow{\ \alpha^\ast\ }T
 \xrightarrow{\ \alpha\ }\mathbb G_{m,S}
\]
be a diagram of groups, where \(T\) is a torus,\footnote{This
qualification is supplied by the editors.} such that
\(\alpha\circ\alpha^\ast=2\).  Let \(R\) be the maximal torus of
\(\Ker(\alpha)\), and put \(K=\alpha^{\ast-1}(R)\).  Then \(K\) is a
subgroup of multiplicative type of \(\mathbb G_{m,S}\); since
\(\alpha\circ\alpha^\ast=2\), it is even a subgroup of
\(\mu_{2,S}\).  In particular,
\[
 K\longrightarrow\mathrm{SL}_L,
 \qquad
 z\longmapsto\begin{pmatrix}z&0\\0&z^{-1}\end{pmatrix}
\]
is central.  Hence there is a central group monomorphism
\[
 K\longrightarrow R\times\mathrm{SL}_L,
 \qquad
 z\longmapsto
 \left(
 \alpha^\ast(z),
 \begin{pmatrix}z&0\\0&z^{-1}\end{pmatrix}
 \right).
\]

Let
\[
 G=(R\times\mathrm{SL}_L)/K.
\]
This is an affine smooth group over \(S\), with connected fibers.
The sequence
\[
 1\longrightarrow K\longrightarrow
 R\times t_S(\mathbb G_{m,S})
 \xrightarrow{\ u\ }T\longrightarrow1,
 \qquad
 u(x,t_S(z))=x\alpha^\ast(z),
\]
is exact.  The image of \(R\times t_S(\mathbb G_{m,S})\) in \(G\)
is therefore a torus \(T'\) isomorphic to \(T\).  If \(\alpha'\) is
the character of \(T'\) transported from \(\alpha\), one readily
checks that \((G,T',\alpha')\) is an elementary \(S\)-system, that
\(\mathfrak g^{\alpha'}\) is isomorphic to \(L\), and that
\(\alpha'^\ast\) is obtained from \(\alpha^\ast\) under the
isomorphism \(T\xrightarrow{\sim}T'\).  Thus the object
\[
 \left(
 \mathbb G_{m,S}\xrightarrow{\ \alpha'^\ast\ }T'
 \xrightarrow{\ \alpha'\ }\mathbb G_{m,S},\,
 \mathfrak g^{\alpha'}
 \right)
\]
of the category \(\mathcal D\) defined in \SGARefLink{sga3:xx:4.2}{4.2} is isomorphic to
\[
 \left(
 \mathbb G_{m,S}\xrightarrow{\ \alpha^\ast\ }T
 \xrightarrow{\ \alpha\ }\mathbb G_{m,S},\,L
 \right).
 \]
We have therefore proved:

\paragraph{\textbf{Theorem 5.11.}}
\SGARefTarget{sga3:xx:5.11}
\emph{With the notation of \SGARefLink{sga3:xx:4.2}{4.2}, the functor}
\[
 (G,T,\alpha)\longmapsto
 \left(
 \mathbb G_{m,S}\xrightarrow{\ \alpha^\ast\ }T
 \xrightarrow{\ \alpha\ }\mathbb G_{m,S},\,
 \mathfrak g^\alpha
 \right)
\]
\emph{is an equivalence of categories between
\(\mathcal E\) and \(\mathcal D\).}

